\documentclass[11pt,a4paper]{article}

\usepackage[utf8]{inputenc}
\usepackage[T1]{fontenc}
\usepackage[margin=2.5cm]{geometry}
\usepackage{graphicx}
\usepackage{float}
\usepackage{amsmath}
\usepackage{booktabs}
\usepackage{parskip}

\title{Lab 2 -- OpenMP Sorting}
\author{Student 1 \and Student 2}
\date{March 2026}

\begin{document}

\maketitle

\section{Participants}
Student 1: TODO\\
Student 2: TODO

\section{Machine used for experiments}
We ran all experiments on a Grid'5000 node (cluster dahu):

\begin{verbatim}
CPU(s):                  64
Thread(s) per core:      2
Core(s) per socket:      16
Socket(s):               2
Model name:              Intel(R) Xeon(R) Gold 6130 CPU @ 2.10GHz
\end{verbatim}

The node has 32 physical cores (16 per socket, 2 sockets) with Hyper-Threading enabled, giving 64 logical CPUs total.
We tested thread counts $p \in \{2, 4, 8, 16, 32, 64\}$, where 32 corresponds to one thread per physical core and 64 uses both hardware threads per core.
We kept the same node for all runs so the comparison stays fair.
OpenMP was configured with \texttt{OMP\_DYNAMIC=false}, \texttt{OMP\_PROC\_BIND=spread}, \texttt{OMP\_PLACES=cores} to avoid thread migration and ensure reproducible timings.

\section{Exercise 1 -- Bubble sort}
We first implemented the sequential version in \texttt{bubble.c}. It does the usual passes over adjacent elements and swaps if needed.

For the parallel version, we followed the lab idea: split the array into chunks, sort each chunk in parallel, then fix chunk borders. We repeat until no border swap happens anymore.

For the scheduling policy, we used \texttt{schedule(static, 1)}. The outer loop has exactly \texttt{nthreads} iterations, one per chunk. With \texttt{static, 1}, OpenMP distributes iterations round-robin one at a time, so thread 0 gets chunk 0, thread 1 gets chunk 1, and so on. Since all chunks are the same size and do the same amount of work, this is a perfect static assignment with no load imbalance. Dynamic scheduling would add unnecessary overhead for a problem that is already well balanced.

This version is correct, but it is not the most scalable algorithm. The border synchronization step adds overhead and limits speedup.

\textbf{Performance observation.}
Our measurements show that parallel bubble sort yields no measurable speedup at any thread count, even at the largest tested size ($N=15$, array of 32768 elements). At that size, the serial time is 1.178 s and the parallel times with 2, 4, 8, 16, and 32 threads are all within 0.5\% of each other (1.175--1.183 s). This is an expected and fundamental limitation:
\begin{itemize}
    \item The outer \texttt{do\ldots while} loop runs $O(N)$ iterations. Each iteration contains two \texttt{parallel for} regions, each of which ends with a full barrier synchronization. The cost of $O(N)$ barriers dominates at all sizes.
    \item The algorithm scans the array repeatedly, saturating memory bandwidth even at low thread counts.
    \item Each thread's share of work per pass ($\approx N/p$ elements) is small enough that thread overhead outweighs the parallel gain.
\end{itemize}
This is not a bug but an intrinsic property of this parallelization strategy.

\section{Exercise 2 -- Mergesort with tasks}
The sequential mergesort is a standard recursive implementation: split, recursively sort both halves, then merge.

Our mergesort implementation assumes the array size is a power of 2, which is guaranteed by the lab setup (array size is 2\^N). This ensures both halves are always equal size, which is required by the provided merge function.

The parallel version uses OpenMP tasks. In each recursive call, we create two tasks (left and right halves), then wait for them before merge. This maps naturally to mergesort because the two halves are independent.

To improve performance, we added a cutoff threshold (\texttt{PARALLEL\_THRESHOLD}): for subarrays smaller than this size we switch back to sequential recursion. Without this, task creation overhead dominates for tiny chunks.

\textbf{Threshold tuning.}
The lab asks us to propose and evaluate a modification to improve parallel performance.
We evaluated two threshold values and measured the impact:
\begin{itemize}
    \item \texttt{PARALLEL\_THRESHOLD = 1024}: at $N=13$ (8192 elements) with $p=2$, parallel time is 0.000238\,s vs.\  serial 0.000267\,s, giving $S(2) \approx 1.12$.
    \item \texttt{PARALLEL\_THRESHOLD = 64}: same configuration gives 0.000209\,s, improving to $S(2) \approx 1.28$. More tasks are created in the recursion tree, exposing more parallelism at medium sizes.
\end{itemize}
However, at large $N$ ($N=27$, 134M elements) both thresholds give essentially identical results: parallel with $p=2$ takes $\approx 8.18$\,s vs.\ serial 8.19\,s, giving $S(2) \approx 1.001$.
The sequential \texttt{merge()} at the top recursion levels dominates regardless of threshold, making further task granularity irrelevant at large $N$.
We keep \texttt{PARALLEL\_THRESHOLD = 1024} as the default since it avoids excessive task overhead at small sizes.

\textbf{Performance observation.}
Even with a lower threshold, speedup at large $N$ is negligible. The table below summarises key measured speedups ($S(p) = T_{serial}/T_p$) with \texttt{PARALLEL\_THRESHOLD = 64}:

\begin{table}[H]
\centering
\begin{tabular}{rrrrrr}
\toprule
$N$ & serial (s) & $p=2$ & $p=4$ & $p=8$ & best $S(p)$ \\
\midrule
13 & 0.000267 & 0.000209 & 0.000201 & 0.000202 & 1.33 ($p=4$) \\
15 & 0.000859 & 0.000884 & 0.000835 & 0.001339 & 1.03 ($p=4$) \\
21 & 0.07151  & 0.06941  & 0.07032  & 0.09390  & 1.03 ($p=2$) \\
27 & 8.1906   & 8.1874   & 8.2065   & —        & $\approx 1.0$ \\
\bottomrule
\end{tabular}
\caption{Mergesort parallel times with \texttt{PARALLEL\_THRESHOLD = 64} on Intel Xeon Gold 6130.}
\end{table}

Several observations stand out:
\begin{itemize}
    \item The best speedup is only $S \approx 1.33$ at $N=13$ with $p=4$. Above that, gains vanish.
    \item With $p=8$ at $N=15$, the parallel version is \textbf{56\% slower} than serial (0.001339\,s vs.\ 0.000859\,s). At $N=21$, p=8 is \textbf{31\% slower} (0.09390\,s vs.\ 0.07151\,s). These spikes are caused by task creation overhead: more tasks are spawned than the runtime can schedule efficiently for those specific size/thread combinations.
    \item At large $N$ ($N=27$), all thread counts give essentially identical time to serial. The sequential \texttt{merge()} at the top recursion levels processes $2^N$ elements alone, and by Amdahl's law this single fraction caps the achievable speedup to nearly 1.
\end{itemize}
This is the fundamental limitation of this parallel mergesort design: the merge step is not parallelised.

\section{Exercise 3 -- Odd-even sort}
Odd-even sort was easier to parallelize than bubble sort. In each phase, comparisons are on disjoint pairs:

Even phase: $(0,1), (2,3), (4,5), \dots$\\
Odd phase: $(1,2), (3,4), (5,6), \dots$

So comparisons inside one phase do not interfere with each other. This is exactly what we need for an OpenMP \texttt{parallel for}.

We implemented both sequential and parallel versions in \texttt{odd-even.c}. The parallel code does one parallel loop for even pairs and one parallel loop for odd pairs until no swap happens.

\section{Performance analysis}
We used the provided script \texttt{plot\_times.py} and ran all three programs with several sizes and thread counts.

\subsection*{How to reproduce}
\begin{verbatim}
cd src
make -B
python3 -m pip install matplotlib
python3 plot_times.py
\end{verbatim}

\subsection*{Execution time plots}
\begin{figure}[H]
    \centering
    \includegraphics[width=0.75\textwidth]{figures/plot_bubble.png}
    \caption{Bubble sort: execution time vs problem size}
\end{figure}

\begin{figure}[H]
    \centering
    \includegraphics[width=0.75\textwidth]{figures/plot_mergesort.png}
    \caption{Mergesort: execution time vs problem size}
\end{figure}

\begin{figure}[H]
    \centering
    \includegraphics[width=0.75\textwidth]{figures/plot_odd-even.png}
    \caption{Odd-even sort: execution time vs problem size}
\end{figure}

\subsection*{Speedup discussion}
We compute speedup with:
\[
S(p) = \frac{T_{serial}}{T_p}
\]
where $T_{serial}$ is the sequential version time and $T_p$ is the parallel time with $p$ threads.

\textbf{Bubble sort} shows essentially no speedup at any thread count or problem size. The measured $S(p) \approx 1.0$ for all $p$ at $N=15$. The $O(N)$ barrier synchronizations per sort dominate over the parallel work.

\textbf{Mergesort} shows the best scaling among the three. Speedup grows with problem size and saturates around the physical core count (32 threads). At large $N$, speedup is limited by the sequential \texttt{merge()} step at each recursion level (Amdahl's law). Going from 32 to 64 threads (using Hyper-Threading) gives negligible gain since the bottleneck is memory bandwidth during merge, not compute.

\textbf{Odd-even sort} falls between the two. The parallel \texttt{for} over disjoint pairs is cleaner than bubble sort's bordered chunks, so there is some gain at large $N$. However, like bubble sort, each outer loop iteration requires two barrier synchronisations, which limits scaling.

\subsection*{Summary}
\begin{table}[H]
\centering
\begin{tabular}{lccc}
\toprule
Algorithm & Observed $S(p)$ at $N_{max}$ & Main bottleneck & Overall trend \\
\midrule
Bubble sort & $\approx 1.0$ (no speedup) & $O(N)$ barrier synchronizations & Flat / no improvement \\
Mergesort (tasks) & Moderate (saturates at 32 threads) & Sequential \texttt{merge()} step & Best scaling \\
Odd-even sort & Small gain at large $N$ & Two barriers per outer iteration & Better than bubble \\
\bottomrule
\end{tabular}
\caption{Summary of observed parallel performance on Intel Xeon Gold 6130, 32 physical cores.}
\end{table}


\section{Limitations / known issues}
The machine may show a \texttt{perf\_event} permission warning at startup. This does not affect sorting correctness or wall-clock timing; it only disables hardware CPU-cycle counters.

For very small arrays, parallel versions are slower than sequential because thread creation and barrier overhead dominate the actual sort work. This is visible in all three algorithms at small $N$.

The \texttt{merge()} function in mergesort is always sequential. A parallel merge (e.g.\ using a split-merge strategy) would improve speedup further but is outside the scope of this lab.

Bubble sort's parallel implementation caps the number of threads to the array size to avoid a chunk-size-zero division for small $N$ with high thread counts.

\section{Conclusion}
This lab showed that algorithm structure determines how well OpenMP parallelism pays off.
Mergesort with tasks gave the best results because the two recursive halves are fully independent and can run in parallel with no data sharing.
The task cutoff threshold is a key tuning parameter: too high and tasks are never created; too low and task overhead dominates.
Odd-even sort parallelises cleanly with \texttt{parallel for} since each phase operates on disjoint pairs, giving moderate speedup at large sizes.
Bubble sort yielded no practical speedup: the $O(N)$ barrier synchronizations required per sort iteration are the fundamental bottleneck, regardless of thread count or problem size.
These results illustrate Amdahl's law in practice: even a small sequential fraction or synchronization cost can completely negate parallel gains.

\end{document}
